%
% 12-geometrisch.tex -- Die geometrische Reihe
%
% (c) 2026 Prof Dr Andreas Müller
%
\subsection{Die geometrische Reihe
\label{buch:analytisch:potenzreihen:subsection:geometrisch}}
Die geometrische Reihe ist die Reihe
\index{geometrische Reihe}%
\[
a+aq+aq^2 + \dots = \sum_{k=0}^\infty aq^k.
\]
Es ist wohlbekannt, dass diese Reihe für $q<1$ konvergent ist und
die Summe
\[
\sum_{k=0}^n aq^k
=
a\frac{q^{n+1}-1}{q-1}
\qquad\Rightarrow\qquad
\sum_{k=0}^n aq^k
=
\lim_{n\to\infty}
a\frac{q^{n+1}-1}{q-1}
=
a\frac{1}{1-q}
\]
hat.

Für komplexe Zahlen $z\in\mathbb{C}$ mit $|z|<1$ definiert die
geometrische Reihe die holomorphe Funktion
\[
\sum_{k=1} z^k
=
\frac{1}{1-z}
=
f(z).
\]
Der Konvergenzradius der Reihe ist $\varrho=1$.
\index{Konvergenzradius}%
Die Polstelle bei $z=1$ führt zu einer divergenten Potenzreihe.
\index{Polstelle}%
\index{divergent}%

Die Funktion 
\begin{equation}
z
\mapsto
\frac{1}{x-z\mathstrut}
\label{buch:analytisch:eqn:1xz}
\end{equation}
ist ebenfalls holomorph, wir möchten sie jetzt in eine im Punkt $a$
zentrierte Potenzreihe entwickeln.
Dazu verwenden wir die Identität
\begin{equation}
\frac{1}{x-z\mathstrut}
=
\frac{1}{(x-a)-(z-a)\mathstrut}
=
\frac{1}{\displaystyle(x-a)\biggl(1-\frac{z-a\mathstrut}{x-a\mathstrut}\biggr)}
=
\frac{1}{x-a}
\cdot
\frac{1}{\displaystyle1-\frac{z-a\mathstrut}{x-a\mathstrut}}.
\label{buch:analytisch:eqn:geometrisch1}
\end{equation}
Der letzte Faktor auf der rechten Seite lässt sich mit der geometrischen
Reihe als Potenzreihe der Form
\begin{equation}
\frac{1}{\displaystyle1-\frac{z-a\mathstrut}{x-a\mathstrut}}
=
\sum_{k=0}^\infty \biggl(\frac{z-a\mathstrut}{x-a\mathstrut}\biggr)^k
=
\sum_{k=0}^\infty
c_k
(z-a)^k
\qquad\text{mit $c_k=1/(x-a)^k$}
\label{buch:analytisch:eqn:geometrischintegriert}
\end{equation}
entwickeln.
